% TOP_PROBLEM: {"id": 2307032, "title": "Research Problems in Function Theory — Problem 7.32", "category_id": 8, "category": {"id": 8, "name": "analysis", "display_name": "Analysis", "description": "Limits, continuity, calculus, and function theory.", "slug": "analysis", "order_index": 8, "created_at": "2026-07-31T15:26:25.671Z"}, "status": "open"}
% FIRST_POSTED: null
% ENTRY_KIND: statement_only
% Imported from: batch05.tex; UnsolvedMath ID 2307032
\documentclass[11pt]{article}
\usepackage{amsmath,amssymb,mathtools}
\usepackage{fontspec,unicode-math}
\setmainfont{Latin Modern Roman}
\setmathfont{Latin Modern Math}
\usepackage{xurl,hyperref}
\newcommand{\sref}[2]{\hyperlink{source:#1:#2}{[S#2]}}
\newcommand{\cataloguescope}[1]{\paragraph{Mathematical statement.}#1}
\begin{document}
% UnsolvedMath ID 2307032; source catalogue code AMR-022-7032
\setcounter{section}{2307031}
\section{Research Problems in Function Theory — Problem 7.32}\label{2307032}
{\small\sffamily UnsolvedMath ID 2307032 · Analysis}
\subsection{Definitions and mathematical statement}

\paragraph{Shared notation.}$\mathbb N=\{1,2,\ldots\}$, and $\mathbb Z,\mathbb Q,\mathbb R,\mathbb C$ denote the integers, rationals, reals and complex numbers. Any other conventions are specified locally.


For a continuous real function $\mu$ on $[0,1]$, define
\[\omega_1(h,\mu)=\sup_{|s-t|\le h,\ s,t\in[0,1]}|\mu(s)-\mu(t)|,\]
\[\omega_2(h,\mu)=\sup_{0\le s\le h,\ t,t+s,t+2s\in[0,1]}|\mu(t+2s)-2\mu(t+s)+\mu(t)|.\]
A singular increasing function here is nonconstant, continuous and nondecreasing, with derivative zero Lebesgue-almost everywhere. Big-$O$ constants may depend on the chosen functions but not on $h$.

\paragraph{Statement recorded in the supplied source.}
For every function $\phi:(0,1]\to(0,\infty)$ increasing to $+\infty$ as its argument decreases to zero, does there exist a singular increasing $\mu$ for which both
\[\omega_1(h,\mu)=O(h\phi(h)),\qquad
\omega_2(h,\mu)=O\!\left(h(\log(1/h))^{-1/2}\right)\quad(h\downarrow0)
\]
hold? \sref{2307032}{2}
The update reports an affirmative answer by Anderson, Fernández and Shields. \sref{2307032}{2}
\subsection{Short English statement}
Can the two sharp borderline bounds on continuity and second differences be attained simultaneously by a singular increasing function, for every prescribed slowly diverging factor?
\subsection{Sources}
\begin{description}
\item[\hypertarget{source:2307032:1}{S1}] ULAM.AI and the UnsolvedMath Contributors, \emph{UnsolvedMath: A Curated Collection of Open Mathematics Problems}, record 2307032 (AMR-022-7032). CC BY4.0. \url{https://huggingface.co/datasets/ulamai/UnsolvedMath}.
\item[\hypertarget{source:2307032:2}{S2}] W. K. Hayman and E. F. Lingham, \emph{Research Problems in Function Theory} (2018), Chapter7, Problem7.32 and its complete update. \url{https://arxiv.org/abs/1809.07200}.
\end{description}
\label{end:2307032}
\end{document}
